\documentclass[11pt]{article}
\usepackage[T1]{fontenc}
\usepackage[utf8]{inputenc}
\usepackage{lmodern}
\usepackage[margin=1in]{geometry}
\usepackage{microtype}
\usepackage{xcolor}
\usepackage{amsmath,amssymb,amsthm}
\usepackage{booktabs,tabularx,array}
\usepackage{enumitem}
\usepackage{hyperref}
\usepackage{titlesec}
\usepackage{fancyhdr}
\usepackage{lastpage}

\definecolor{TitleBlue}{HTML}{1F4E79}
\definecolor{RuleGray}{HTML}{D9E1E8}
\definecolor{SoftGray}{HTML}{4A4A4A}
\definecolor{LightBlue}{HTML}{EEF5FB}
\definecolor{BorderBlue}{HTML}{A9C2D9}

\hypersetup{
  colorlinks=true,
  linkcolor=TitleBlue,
  urlcolor=TitleBlue,
  citecolor=TitleBlue,
  pdftitle={Technical note on Closing the Extension Gap: Decidability of ALCIRCC5 via Direct Model Construction},
  pdfauthor={OpenAI}
}

\titleformat{\section}{\large\bfseries\color{TitleBlue}}{\thesection}{0.6em}{}
\titleformat{\subsection}{\normalsize\bfseries\color{TitleBlue}}{\thesubsection}{0.6em}{}
\titlespacing*{\section}{0pt}{1.2em}{0.55em}
\titlespacing*{\subsection}{0pt}{0.9em}{0.4em}

\setlength{\parindent}{0pt}
\setlength{\parskip}{0.6em}
\renewcommand{\arraystretch}{1.18}
\newcolumntype{Y}{>{\raggedright\arraybackslash}X}
\newcolumntype{C}{>{\centering\arraybackslash}m{1.9cm}}

\pagestyle{fancy}
\fancyhf{}
\fancyhead[L]{\footnotesize Technical note on \emph{Closing the Extension Gap}}
\fancyhead[R]{\footnotesize Page \thepage\ of \pageref{LastPage}}
\renewcommand{\headrulewidth}{0.25pt}
\renewcommand{\footrulewidth}{0pt}

\newcommand{\DR}{\mathsf{DR}}
\newcommand{\PO}{\mathsf{PO}}
\newcommand{\PP}{\mathsf{PP}}
\newcommand{\PPI}{\mathsf{PPI}}
\newcommand{\EQ}{\mathsf{EQ}}
\newcommand{\comp}{\mathsf{comp}}
\newcommand{\Llab}{\mathcal{L}}
\newcommand{\Elab}{\mathcal{E}}
\newcommand{\DNsafe}{\mathsf{DN}_{\mathsf{safe}}}

\begin{document}

{\centering
  {\LARGE\bfseries\color{TitleBlue} Technical note on the manuscript}\par
  \vspace{0.3em}
  {\large\bfseries \emph{Closing the Extension Gap: Decidability of ALCIRCC5 via Direct Model Construction}}\par
  \vspace{0.5em}
  {\normalsize A concise written version of the assessment discussed in chat}\par
}

\vspace{1em}
\noindent\fcolorbox{BorderBlue}{LightBlue}{%
\parbox{0.965\linewidth}{%
\textbf{Bottom line.}
I do not agree that the manuscript proves decidability of $\mathcal{ALCI}_{RCC5}$.
There is a useful local observation in the paper - the $\DR/\PPI$ configuration in Lemma 4.1 really does force a $\DR$ edge by composition filtering - but the headline theorem still fails for two independent reasons:
\begin{enumerate}[leftmargin=1.5em,itemsep=0.2em,topsep=0.45em]
  \item the algebraic analysis in Lemma 3.2 and Remark 4.3 misreads the RCC5 table, and
  \item the key new step, Theorem 5.5 on path-consistency of the disjunctive network, is not proved and is in fact false for the network defined in Definition 5.3.
\end{enumerate}
}}

\section{Context and verdict}

The manuscript claims to close the open decidability problem for $ALCIRCC5$ by supplying the missing soundness step for the tableau construction \cite{Claude2026}. That is an important target: Wessel's report explicitly leaves $ALCIRCC5$ and $ALCIRCC8$ open, while giving only lower complexity bounds and noting that the missing ingredient is a convincing infinity-handling or blocking style argument for tableau-based reasoning \cite{Wessel2003}.

My assessment is narrower. The paper contains a worthwhile local diagnosis - especially the observation that the tableau's own composition filtering already forces some of the edge adjustments one would want in an unraveling argument - but the manuscript does not yet deliver a correct proof of decidability.

\section{First problem: the algebraic analysis is inconsistent}

The RCC5 table used in the manuscript is indexed so that the entry in row $S(b,c)$ and column $R(a,b)$ is $\comp(R,S)$. With that indexing, Table~1 gives
\[
\comp(\DR,\PP)=\{\DR,\PO,\PP\},
\qquad
\comp(\PP,\DR)=\{\DR\}.
\]
But Lemma 3.2 states instead that
\[
\comp(\DR,\PP)=\{\DR\},
\]
which is false. The proof table printed immediately below the lemma repeats the same mistake.

This matters again in Remark 4.3. The remark wants a dual to the true statement
\[
\PPI \notin \comp(\DR,\PPI)=\{\DR\},
\]
and then informally treats the ``dual'' $\DR/\PP$ configuration as if it were governed by the same singleton composition value. The correct dual-looking statement is
\[
\PP \notin \comp(\PP,\DR)=\{\DR\},
\]
not the false statement with arguments reversed. So the manuscript blends three different claims:
\begin{align*}
\PPI &\notin \comp(\DR,\PPI)=\{\DR\}, \\
\PP &\notin \comp(\PP,\DR)=\{\DR\}, \\
\PP &\notin \comp(\DR,\PP),
\end{align*}
where only the first two are true.

This does not destroy Lemma 4.1 itself; that local lemma is correct. But it does show that the higher-level algebraic narrative in Section 3 is not reliable as written.

\section{Second problem: Theorem 5.5 is not established}

The manuscript's new global step is Definition 5.3 plus Theorem 5.5. For non-parent-child pairs in the unraveling, the domain of admissible relations is defined only from endpoint types:
\[
D(d_1,d_2)=\DNsafe(\operatorname{tp}(d_1),\operatorname{tp}(d_2)).
\]
So all information about which tableau representatives actually realized those types is discarded.

That is exactly where the proof of Theorem 5.5 breaks. In Case B, the argument starts with an arbitrary
\[
R_{12}\in D(d_1,d_2),
\]
so there are tableau nodes $a_1,a_2$ of the correct types with
\[
\Elab(a_1,a_2)=R_{12}.
\]
The proof then introduces a witness node $w'$ coming from the actual tableau edge realizing the fixed tree relation
\[
R_{23}=S.
\]
However, nothing guarantees that $a_2=T(d_2)$. If $a_2\neq T(d_2)$, then one only knows
\[
\Elab(a_2,w')\in \text{some RCC5 value compatible with the older graph},
\]
not that $\Elab(a_2,w')=S$. Consequently composition filtering yields only
\[
\Elab(a_1,w')\in \comp\bigl(R_{12},\Elab(a_2,w')\bigr),
\]
not the required statement
\[
\Elab(a_1,w')\in \comp(R_{12},S).
\]
This is the missing step. The manuscript notices the issue informally in Section 7.4, where it says Cases B and C ``could benefit from a more explicit construction.'' In my view that caveat understates the problem: Cases B and C are precisely where the proof fails.

\section{A concrete counterexample to Theorem 5.5}

The failure is not just a proof gap. The network of Definition 5.3 is not path-consistent in general.

Consider the concepts
\[
B := \exists \DR.(p \sqcup q),
\qquad
A := \exists \DR.B \sqcap \exists \PO.B,
\qquad
C_0 := \exists \DR.A.
\]
Take a saturated completion graph with node labels
\[
\Llab(x)=\{C_0\}, \quad \Llab(a)=\{A\}, \quad \Llab(b_1)=\Llab(b_2)=\{B\},
\quad \Llab(c_1)=\{p\}, \quad \Llab(c_2)=\{q\},
\]
and with the following RCC5 edges on distinct nodes:

\begin{center}
\begin{tabularx}{0.96\linewidth}{@{}>{\bfseries}lYY@{}}
\toprule
Required witnesses & $x\,\DR\,a$, $a\,\DR\,b_1$, $a\,\PO\,b_2$, $b_1\,\DR\,c_1$, $b_2\,\DR\,c_2$ \\
Type-coarse duplicates & $x\,\DR\,b_1$, $x\,\PO\,b_2$ \\
Remaining edges & $x\,\PP\,c_1$, $x\,\DR\,c_2$, $a\,\DR\,c_1$, $a\,\DR\,c_2$, $b_1\,\DR\,b_2$, $b_1\,\DR\,c_2$, $b_2\,\PO\,c_1$, $c_1\,\DR\,c_2$ \\
\bottomrule
\end{tabularx}
\end{center}

This atomic RCC5 network is consistent; one can check that every triple satisfies the RCC5 composition table. Also, $b_2$ is blocked by $b_1$ under the manuscript's blocking test, since they have the same label $\{B\}$.

Now move to the tree unraveling and consider:
\begin{itemize}[leftmargin=1.5em]
  \item $d_1$ = the root copy of $x$,
  \item $d_2$ = the copy of the active $B$-node $b_1$ introduced under $a$,
  \item $d_3$ = the $\DR$-child copy of $c_1$ under $d_2$.
\end{itemize}

By Definition 5.3 of the manuscript,
\[
D(d_1,d_2)=\DNsafe(\operatorname{tp}(d_1),\operatorname{tp}(d_2))=\{\DR,\PO\},
\]
because in the completion graph $x$ has both a $\DR$-edge to $b_1$ and a $\PO$-edge to $b_2$, and $b_1,b_2$ have the same type $B$.

The tree edge from $d_2$ to $d_3$ is fixed, so
\[
R_{23}=\DR.
\]
Also,
\[
D(d_1,d_3)=\DNsafe(\operatorname{tp}(d_1),\operatorname{tp}(d_3))=\{\PP\},
\]
because $c_1$ is the unique node of its type and the completion graph contains the edge $x\,\PP\,c_1$.

Now choose
\[
R_{12}=\PO\in D(d_1,d_2).
\]
Then the RCC5 table gives
\[
\comp(\PO,\DR)=\{\DR,\PO,\PPI\}.
\]
Therefore
\[
\comp(\PO,\DR)\cap D(d_1,d_3)
=
\{\DR,\PO,\PPI\}\cap\{\PP\}
=
\varnothing.
\]
So the network $\mathcal N$ from Definition 5.3 is \emph{not} path-consistent. This directly contradicts Theorem 5.5.

\section{What survives, and what would need to change}

The paper does isolate a real local phenomenon: the configuration
\[
\Elab(n,w)=\DR, \qquad \Elab(n,m)=\PPI
\]
does force $\Elab(w,m)=\DR$ by composition filtering. That is a sound and useful observation. But it is not enough to support the global model construction built in Section 5.

The reason is structural: the domains $D(d_1,d_2)$ depend only on endpoint types. Once edges are quotiented that coarsely, the path-consistency argument loses the very information it needs in Cases B and C, namely which concrete predecessor context realized a given relation. The counterexample above shows that endpoint-type domains can accumulate incompatible relations from distinct tableau representatives of the same type.

A plausible repair would have to refine the edge domains by richer blocker context than bare node type - for example by storing predecessor-color information, witness profile information, or another finite signature strong enough to guarantee replay of the fixed tree edge. But that is a different theorem from the one currently proved in the manuscript.

\section{Conclusion}

My conclusion is therefore:
\begin{itemize}[leftmargin=1.5em]
  \item I agree with the local forcing observation in Lemma 4.1.
  \item I do \emph{not} agree that the manuscript proves decidability of $\mathcal{ALCI}_{RCC5}$.
  \item The decisive obstacle is Theorem 5.5: the proof does not go through, and the theorem is false for the network defined in Definition 5.3.
\end{itemize}

So the manuscript is a useful discussion piece, but it does not yet close Wessel's original open problem.

\vspace{0.7em}
\noindent\rule{\linewidth}{0.4pt}

\begin{thebibliography}{9}
\bibitem{Claude2026}
Claude.
\newblock \emph{Closing the Extension Gap: Decidability of ALCIRCC5 via Direct Model Construction}.
\newblock Discussion manuscript, April 2026.

\bibitem{Wessel2003}
Michael Wessel.
\newblock \emph{Qualitative Spatial Reasoning with the ALCIRCC Family - First Results and Unanswered Questions}.
\newblock Technical Report FBI-HH-M-324/03, University of Hamburg, 2002/2003.
\end{thebibliography}

\end{document}
